\section{Prise en main d'un modèle de processeur}

\subsection{RISC-V RV32I et chargement des instructions}

Dans un premier temps, nous avons étudié le jeu d'instructions RISC-V RV32I et chargé les instructions dans la mémoire du processeur. Un \textit{testbench} a été créé afin de simuler le processeur et de vérifier le bon fonctionnement des instructions. 
\\

Le processeur n'exécute pas directement les instructions assembleur lisibles par un humain, comme \texttt{addi} ou \texttt{add} : chaque instruction doit être encodée en code machine selon le format défini par le jeu d'instructions. Ces mots binaires de 32 bits sont représentés en hexadécimal. Chaque groupe de huit chiffres hexadécimaux correspond ainsi à une instruction.
\\

L'exemple ci-dessous illustre le chargement d'un programme de test :
\\

\begin{lstlisting}[caption={Exemple de chargement d'instructions}, label={lst:load_instructions}]
// instructions at address 0x0
@0
00200093 // addi x1, x0, 2
00400113 // addi x2, x0, 4
00110133 // add x2, x2, x1
001150b3 // srl x1, x2, x1
0420a023 // sw x2, 64(x1)
00100073 // ebreak
\end{lstlisting}

Les chronogrammes ci-dessous mettent en évidence l'exécution séquentielle de ces instructions, avec le comportement des registres (Figure~\ref{fig:processor_2}) et l'écriture en mémoire (Figure~\ref{fig:processor_3}) :
\begin{itemize}
    \item \textcolor{red}{\textbf{Étape 1}} : le registre \texttt{x1} est initialisé à 2, puis \texttt{x2} à 4.
    \item \textcolor{cyan}{\textbf{Étape 2}} : leur addition est stockée dans \texttt{x2}.
    \item \textcolor{magenta}{\textbf{Étape 3}} : un décalage logique à droite de \texttt{x2} par \texttt{x1} est affecté à \texttt{x1}.
    \item \textcolor{yellow}{\textbf{Étape 4}} : la valeur de \texttt{x2} est sauvegardée en mémoire à l’adresse calculée à partir de \texttt{x1} et de l’offset 64.
\end{itemize}

Les figures suivantes détaillent respectivement le comportement interne des registres (Figure~\ref{fig:processor_2}) et l'écriture effective des données dans l'espace mémoire.

\begin{figure}[H]
    \centering
    \includegraphics[width=0.99\textwidth]{3_INSTRU_0_modified.png}
    \caption{Détail des opérations dans les registres}
    \label{fig:processor_2}
\end{figure}

\begin{figure}[H]
    \centering
    \includegraphics[width=0.99\textwidth]{3_INSTRU_1_modified.png}
    \caption{Validation du stockage de la valeur en mémoire}
    \label{fig:processor_3}
\end{figure}

\subsection{Analyse des modes d'adressage et extension RV32IM}

\subsubsection{Modes d'adressage}
L'architecture RISC-V manipule différents modes d'adressage pour accéder aux données :
\begin{itemize}
    \item  \textbf{Adressage immédiat :} La valeur de la donnée est directement encodée au sein même de l'instruction (ex: \texttt{addi x1, x0, 2}).
    \item  \textbf{Adressage direct :} L'adresse de la mémoire ciblée est explicitement fournie dans l'instruction.
    \item  \textbf{Adressage indirect (basé sur un registre) :} L'adresse mémoire effective est contenue dans un registre, souvent combinée avec un décalage constant (\textit{offset}), comme dans l'instruction de stockage \texttt{sw x2, 64(x1)}.
\end{itemize}

\subsubsection{Évolution vers le jeu d'instructions RV32IM}
Le processeur fourni gère nativement le jeu d'instructions de base RV32I. Pour le rendre compatible avec l'extension \textbf{M} (\textit{Multiplication and Division}), les modifications matérielles suivantes sont requises :
\begin{itemize}
    \item Intégration d'unités de calcul arithmétique dédiées dans l'ALU pour la multiplication et la division entières (instructions \texttt{mul}, \texttt{div}).
    \item Décodage des nouveaux \textit{opcodes} et \textit{funct3/funct7} spécifiques à l'extension M dans l'unité de contrôle.
    \item Adaptation du programme de test (`instr.mem`) pour valider ces nouvelles opérations arithmétiques.
\end{itemize}

\subsection{Passage du jeu d'instructions RV32I au jeu RV32IM}

Voici les modifications apportées au code Verilog du processeur pour supporter l'extension M :

\begin{lstlisting}[caption={Ajout de l'extension M dans le processeur}, label={lst:processor_extension_m}]
// Add functions RV32IM
`define MUL   {3'b000, 7'b0000001}
`define DIV   {3'b100, 7'b0000001}

// Addictionnal function
`MUL: regs[`RD] <= regs[`RS1] * regs[`RS2];
`DIV: regs[`RD] <= regs[`RS1] / regs[`RS2];
\end{lstlisting}

Nous modifions également le fichier `instr.mem` pour inclure des instructions :
\\

\begin{itemize}
    \item \texttt{div} : x2 = x2 / x1 ( x2 = 4 / 2)
    \item \texttt{mul} : x2 = x2 * x1 ( x2 = 2 * 4)
\end{itemize}

\begin{lstlisting}[caption={Instructions pour l'extension M}, label={lst:processor_extension_m_instr}]
// div
00200093 // addi x1, x0, 2
00400113 // addi x2, x0, 4
02114133 // div x2, x2, x1

// mul
00200093 // addi x1, x0, 2
00400113 // addi x2, x0, 4
02110133 // mul x2, x2, x1
\end{lstlisting}

Voici le chronogramme résultant de l'exécution de ces instructions (Figure~\ref{fig:processor_4}) :

\begin{figure}[H]
    \centering
    \includegraphics[width=0.99\textwidth]{3_division_modified.png}
    \caption{Chronogramme du processeur pour division}
    \label{fig:processor_4}
\end{figure}


\begin{figure}[H]
    \centering
    \includegraphics[width=0.99\textwidth]{3_multiplication_modified.png}
    \caption{Chronogramme du processeur pour multiplication}
    \label{fig:processor_5}
\end{figure}
